% Cote 119: offprints, reports, courses, CV and letters (typed and printed part): archive cote 119, site page 249 (PDF page 250).
% typescript. Draft compared: none.
% AI-transcribed by Claude Opus 5.5 (claude-opus-5-5, adaptive thinking, effort high), Message Batches API, batch pilotnext (msgbatch_01A9sdu8SYq23bEtqX9tUncQ), 2026-11-13; not yet checked by eye.
\scanpage{119/249}
\mspage{9}

conditions favorables à son épanouissement qui ont jusqu'à présent fait
défaut.

En plus des maîtres eux-mêmes, les étudiants ont été évidemment les
principales victimes de la survie d'un système de formation et de recrute-%
ment des maîtres basé presque exclusivement sur le critère de l'aptitude
à la recherche. Il est urgent de mettre à profit dans l'immédiat la force
vive libérée aujourd'hui par le mouvement de revendications étudiantes, et
les prises de conscience qu'elles ont précipité\hand{es} chez les enseignants,
pour mettre au point des mécanismes assez souples pour distinguer et coordon-%
ner toutàlafois les deux principales vocations de l'Université, et capables
de s'adapter au fur et à mesure aux besoins,changeants d'une période à l'autre
et d'une Faculté à l'autre. Ces questions sont l'objet dès à présent de
discussions entre étudiants et enseignants (assistants, maîtres assistants,
maîtres de conférences et professeurs), tout au moins au département de
Mathématiques de la Faculté des Sciences d'Orsay, et sans doute un peu
partout dans les universités en effervescence créatrice. Il faut que nous,
les principaux concernés,étudiants et maîtres enseignants et chercheurs,
fassions tout notre possible pour que des solutions d'ensemble soient déga-%
gées dans le plus proche avenir, et puissent être mises en application dès
la rentrée prochaine. \deleted{Il n'est pas dans mon propos ici de passer}\inserted{Je passerai} en revue
\deleted{les}\inserted{plus bas} diverses modifications en profondeur des structures et des habitudes de
pensée dont la nécessité \inserted{me} semble dès à présent acquise. Certaines parmi les
plus importantes découlent logiquement de l'état de faits que je viens
d'essayer d'esquisser. Cet état de faits est bien familier sans doute à ceux
de mes collègues ayant assez d'expérience de l'enseignement ou de la recher-%
che, mais moins familier déjà, semble-t-il, aux assistants ou aux étudiants,
et \over{à}{\unsure{a}} fortiori aux personnes étrangères à l'Université. Mon but dans \over{l'}{la} \inserted{première partie de cet} exposé
\deleted{qui précède n'a pas été de proposer des \uline{solutions}, laissant ce soin à telles}
\deleted{assemblées spontanées plus autorisées que moi qui se réuniront pour en}
\deleted{dégager en commun, mais de contribuer à mettre} \inserted{a été de mettre} en évidence le plus claire-%
ment \deleted{que je peux}\inserted{possible} un des vices de base du système actuel, et \over{d}{\unsure{d|à}}'inciter tous

\margin{249}
